% !PRINTED 252
\weberpage{252}

\begin{center}
\S.\,79.

\gesperrt{The Bezoutiant and its significance for the reality of the roots.}
\end{center}

For the investigation of the reality of the roots of an arbitrary real equation in more general cases, the function which we designated in \S.\,73 as the Bezoutiant of the equation can be applied with advantage.

If
\weq{1}{f(x) = x^n + a_1 x^{n-1} + a_2 x^{n-2} + \cdots + a_n = 0}
is any real equation of degree $n$, then the Bezoutiant was defined as follows:

One sets
\weq{2}{y = t_{n-1} f_0(x) + t_{n-2} f_1(x) + \cdots + t_0 f_{n-1}(x),}
in which the $t_0, t_1 \ldots t_{n-1}$ denote indeterminate variables and
\[
f_0(x) = 1, \quad f_1(x) = x + a_1, \quad f_2(x) = x^2 + a_1 x + a_2 \ldots
\]
Then [\S.\,73, (10)]
\weq{3}{S(y^2) = \frac{1}{n}\left[S(y)\right]^2 + B,}
where $B$ is a quadratic form of the $n - 1$ variables $t_0, t_1 \ldots t_{n-2}$ with coefficients that are composed rationally from the coefficients $a_i$. This function $B$ we have defined as the Bezoutiant, and we have actually constructed it for the cases $n = 4$ and $n = 5$.

Now in the sum on the left-hand side of (3)
\weq{4}{y_1^2 + y_2^2 + y_3^2 + \cdots + y_n^2,}
the $y_1, y_2 \ldots y_n$ are linear functions of the variables $t_0, t_1 \ldots t_{n-1}$, and thus in the formula (3) we have a transformation of the quadratic form
\weq{5}{\frac{1}{n}\left[S(y)\right]^2 + B = \Phi(t_0, t_1 \ldots t_{n-1})}
\gesperrt{into a sum of $n$ squares.}

For the present we make no restrictive assumptions whatever about equality or distinctness of the roots of $f(x)$, and must now first investigate to what extent the linear functions $y_1, y_2 \ldots y_n$ are independent of one another. In this regard the following theorem holds:
%CONT>>
